\documentclass[11pt]{article}

\usepackage[T1]{fontenc}
\usepackage[utf8]{inputenc}
\usepackage{lmodern}
\usepackage{amsmath,amssymb,amsfonts,amsthm}
\usepackage[margin=1in]{geometry}
\usepackage{microtype}
\usepackage{xcolor}
\usepackage[colorlinks=true,linkcolor=teal,urlcolor=teal]{hyperref}

\newcommand{\F}{\mathcal F}
\newcommand{\T}{\mathcal T_1}
\newcommand{\Tr}{\operatorname{Tr}}
\newcommand{\ketbra}[1]{\lvert #1\rangle\!\langle #1\rvert}
\newcommand{\R}{\mathbb R}

\newtheorem*{lemmaD}{Lemma D.4 (rephrased)}

\title{A Short Guide to Lemma D.4\\
\large Tensorizing the one-mode conic stochastic bound}
\author{}
\date{}

\begin{document}
\maketitle

\begin{abstract}
Lemma D.4 lifts the one-mode stochastic inequality of Lemma D.3 to a
multimode map, even when that map couples the modes.  It does not assert a
full multivariate stochastic order.  It proves exactly what is needed later:
an inequality for products of decreasing photon-number weights.  The proof
turns all but one mode into a weighted environment, applies Lemma D.3, and
then repeats the construction mode by mode.
\end{abstract}

\section{The issue addressed by the lemma}

Let
\[
  \Phi_0=\tau_{q_1}\otimes\cdots\otimes\tau_{q_r}
\]
be a product of thermal states.  A multimode channel $\Gamma$ need not be a
product channel: it may create correlations among the output modes.  Hence a
one-mode result cannot be tensorized merely by multiplying $r$ inequalities.

Lemma D.4 overcomes this by inserting the weights on the untested modes into
a new one-mode completely positive map.  The trace of that map records the
weighted expectation still left to control.  Applying the one-mode result
successively strips off one factor at a time.

\section{Background and notation}

\subsection{Thermal modes and their amplified reference states}

For the number operator $\widehat N$ on one-mode Fock space $\F$, let
\[
  \tau_q=(1-q)\sum_{k\geq0}q^k\ketbra{k},
  \qquad 0<q<1.
\]
Fix an intensity gain $\gamma>1$ and write
\[
  q^{(\gamma)}=1-\frac{1-q}{\gamma}.
\]
For $r$ modes, define
\begin{equation}\label{eq:states}
  \Phi_0=\bigotimes_{\ell=1}^r\tau_{q_\ell},
  \qquad
  \tau^{(\gamma)}
  =\bigotimes_{\ell=1}^r\tau_{q_\ell^{(\gamma)}}.
\end{equation}
The joint number basis is indexed by
$\mathbf k=(k_1,\ldots,k_r)$.

\subsection{Separate covariance in every mode}

Write
\[
  D(z)=\bigotimes_{\ell=1}^r D(z_\ell),
  \qquad z\in\mathbb C^r.
\]
The multimode map $\Gamma$ is displacement covariant with amplitude gain
$\sqrt\gamma$ when
\begin{equation}\label{eq:multicov}
  \Gamma(D(z)XD(z)^*)
  =D(\sqrt\gamma\,z)\Gamma(X)D(\sqrt\gamma\,z)^*.
\end{equation}
It is phase covariant in each mode when it intertwines all independent
rotations
\[
  \bigotimes_{\ell=1}^r e^{\mathrm i\theta_\ell\widehat N_\ell}.
\]
Setting all coordinates of $z$ except $z_j$ equal to zero shows that
\eqref{eq:multicov} contains a one-mode covariance relation for each mode.
This observation is what makes the reduction below work.

\subsection{Product test functions}

The lemma considers
\begin{equation}\label{eq:product-weight}
  W(\widehat{\mathbf N})
  =\prod_{\ell=1}^r w_\ell(\widehat N_\ell),
\end{equation}
where each $w_\ell:\mathbb Z_{\geq0}\to\R_+$ is bounded and decreasing.
Because the factors act on different modes, they commute and the product is
a bounded positive operator.

This restricted class is important.  Ordinary one-dimensional stochastic
orders do not automatically imply an order for every decreasing function of
the vector $(N_1,\ldots,N_r)$ when the output modes may be correlated.  The
factorization in \eqref{eq:product-weight} is precisely what permits a
one-mode-at-a-time proof.

\section{The one-mode input: Lemma D.3}

If $\Lambda$ is a one-mode, completely positive trace-nonincreasing map with
the same displacement and phase covariance, Lemma D.3 states that
\begin{equation}\label{eq:D3}
  \Tr[\Lambda(\tau_q)w(\widehat N)]
  \leq
  \Tr\Lambda(\tau_q)\,
  \Tr[\tau_{q^{(\gamma)}}w(\widehat N)]
\end{equation}
for every bounded nonnegative decreasing $w$.  That lemma uses the external
Gu\c{t}\u{a}--Matsumoto least-noise stochastic-order theorem, together with
the internally proved constant-trace Lemma D.1.  Lemma D.4 introduces no new
external input.

\section{Statement of Lemma D.4}

\begin{lemmaD}
Let $\Gamma$ be an $r$-mode completely positive trace-nonincreasing map,
displacement covariant with gain $\sqrt\gamma$, and phase covariant in every
mode.  For the states in \eqref{eq:states} and every product weight
\eqref{eq:product-weight},
\begin{equation}\label{eq:D4}
  \Tr[\Gamma(\Phi_0)W(\widehat{\mathbf N})]
  \leq
  \Tr\Gamma(\Phi_0)\,
  \Tr[\tau^{(\gamma)}W(\widehat{\mathbf N})].
\end{equation}
\end{lemmaD}

Since both the state and the weight on the right are products,
\begin{equation}\label{eq:right-factor}
  \Tr[\tau^{(\gamma)}W]
  =\prod_{\ell=1}^r
    \Tr[\tau_{q_\ell^{(\gamma)}}w_\ell(\widehat N_\ell)].
\end{equation}

\section{Proof walkthrough}

\subsection*{Step 1: normalize the weights}

By multiplying back their suprema after the proof, it is enough to treat
$0\leq w_\ell\leq1$.  If one of the weights is identically zero, the claim
is immediate.  Otherwise replace $w_\ell$ by
$w_\ell/\lVert w_\ell\rVert_\infty$; both sides of \eqref{eq:D4} acquire the
same factor $\prod_\ell\lVert w_\ell\rVert_\infty$.

The upper bound by one ensures that every product of the weights is an
effect.  Inserting its square root therefore preserves the
trace-nonincreasing property.

\subsection*{Step 2: package the remaining modes into a one-mode map}

The manuscript first treats mode $1$.  Put
\[
  W_{>1}=\prod_{\ell=2}^r w_\ell(\widehat N_\ell),
  \qquad
  \Phi_{0,>1}=\bigotimes_{\ell=2}^r\tau_{q_\ell},
\]
and define
\begin{equation}\label{eq:induced-map}
  \Gamma_1(X)
  =\Tr_{2,\ldots,r}\!\left[
    (I\otimes W_{>1}^{1/2})
    \Gamma(X\otimes\Phi_{0,>1})
    (I\otimes W_{>1}^{1/2})
  \right].
\end{equation}
This is a one-mode map: prepare the other input modes thermally, apply
$\Gamma$, test the other output modes with $W_{>1}$, and discard them.

Every operation in \eqref{eq:induced-map} is completely positive.  For
$X\geq0$,
\begin{align*}
  \Tr\Gamma_1(X)
  &=\Tr\!\left[\Gamma(X\otimes\Phi_{0,>1})
                 (I\otimes W_{>1})\right]\\
  &\leq\Tr\Gamma(X\otimes\Phi_{0,>1})
   \leq\Tr X.
\end{align*}
Thus $\Gamma_1$ is trace nonincreasing.

\subsection*{Step 3: verify that the induced map has the right symmetries}

Apply the multimode covariance relation with displacement vector
$(z,0,\ldots,0)$.  The resulting output displacement acts only on mode $1$,
so it commutes with $I\otimes W_{>1}^{1/2}$ and passes through the partial
trace.  Consequently,
\[
  \Gamma_1(D(z)XD(z)^*)
  =D(\sqrt\gamma\,z)\Gamma_1(X)D(\sqrt\gamma\,z)^*.
\]
The same calculation with a phase rotation in mode $1$ proves phase
covariance.  Therefore $\Gamma_1$ satisfies every hypothesis of Lemma D.3.

\subsection*{Step 4: remove the first weight}

The two traces associated with \eqref{eq:induced-map} are
\begin{align}
  \Tr[\Gamma_1(\tau_{q_1})w_1(\widehat N_1)]
  &=\Tr[\Gamma(\Phi_0)W(\widehat{\mathbf N})],
  \label{eq:left-id}\\
  \Tr\Gamma_1(\tau_{q_1})
  &=\Tr[\Gamma(\Phi_0)W_{>1}].
  \label{eq:mass-id}
\end{align}
Applying \eqref{eq:D3} and substituting these identities gives
\begin{equation}\label{eq:strip-one}
  \Tr[\Gamma(\Phi_0)W]
  \leq
  \Tr[\tau_{q_1^{(\gamma)}}w_1(\widehat N_1)]
  \Tr[\Gamma(\Phi_0)W_{>1}].
\end{equation}

\subsection*{Step 5: iterate}

For complete precision, define
\[
  W_{\geq j}=\prod_{\ell=j}^r w_\ell(\widehat N_\ell),
  \qquad W_{>j}:=W_{\geq j+1},
  \qquad W_{\geq r+1}=I,
\]
with identity factors on the modes not displayed.  The one-mode map used at
stage $j$ is explicitly
\[
  \Gamma_j(X)=\Tr_{\ell\ne j}\!\left[
    W_{>j}^{1/2}\,
    \Gamma\!\left(
      \bigotimes_{\ell<j}\tau_{q_\ell}\otimes X\otimes
      \bigotimes_{\ell>j}\tau_{q_\ell}
    \right)
    W_{>j}^{1/2}
  \right].
\]
Here the partial trace discards every output mode except $j$.  The same
complete-positivity, trace, and covariance checks used for $\Gamma_1$ apply.
The weighted and unweighted traces in Lemma D.3 are respectively the two
terms in
\begin{equation}\label{eq:recursion}
  \Tr[\Gamma(\Phi_0)W_{\geq j}]
  \leq
  a_j\,\Tr[\Gamma(\Phi_0)W_{\geq j+1}],
  \qquad
  a_j:=\Tr[\tau_{q_j^{(\gamma)}}w_j(\widehat N_j)].
\end{equation}

Iterating \eqref{eq:recursion} yields
\[
  \Tr[\Gamma(\Phi_0)W]
  \leq
  \left(\prod_{j=1}^r a_j\right)\Tr\Gamma(\Phi_0).
\]
Equation \eqref{eq:right-factor} identifies the product with
$\Tr[\tau^{(\gamma)}W]$, proving \eqref{eq:D4}.

\section{A two-mode version}

For $r=2$, the proof is especially transparent.  It first gives
\[
  \Tr[\Gamma(\tau_{q_1}\otimes\tau_{q_2})(w_1\otimes w_2)]
  \leq a_1\,
  \Tr[\Gamma(\tau_{q_1}\otimes\tau_{q_2})(I\otimes w_2)].
\]
Applying the same construction to mode $2$ gives
\[
  \Tr[\Gamma(\tau_{q_1}\otimes\tau_{q_2})(I\otimes w_2)]
  \leq a_2\Tr\Gamma(\tau_{q_1}\otimes\tau_{q_2}).
\]
The output may be correlated throughout; the proof never assumes that
$\Gamma$ factors.  It uses only the product form of the input and test
weight, plus separate covariance in each mode.

\section{What the lemma does not claim}

Lemma D.4 proves a conic inequality: the factor
$\Tr\Gamma(\Phi_0)\leq1$ keeps track of subnormalization.  It does not claim
that $\Gamma(\Phi_0)$ is below $\tau^{(\gamma)}$ in operator order, nor that
its full joint photon-number law stochastically dominates the product
geometric law for every multivariate decreasing test.  The product class in
\eqref{eq:product-weight} is sufficient because the fidelity witness in
Lemma D.5 factorizes mode by mode.

\section{Dependencies and role in the paper}

The dependency chain is
\[
  \text{Lemma D.1}
  +\text{external least-noise theorem}
  \longrightarrow\text{Lemma D.3}
  \longrightarrow\text{Lemma D.4}.
\]
Beyond D.3, the proof is internal: it constructs the induced map, checks its
covariance, and iterates.  D.4 supplies the factorized fidelity witness in
D.5, hence Theorem 3.4, and is also applied branchwise in the hybrid converse.
Its prerequisites are completely positive maps, partial trace, product
thermal states, covariance, and Lemma D.3.

\end{document}
